%==============================================================================
%  RELATÓRIO 2  —  Entrega: 13/05/2026
%==============================================================================

\section{Relatório 2: Formulação Matemática e Primeiros Resultados}

\subsection{Introdução}

Este relatório apresenta a segunda etapa do projeto. Com base na concepção
do Relatório 1, são apresentados: a revisão bibliográfica que fundamenta a
técnica, a formulação matemática completa e aprimorada do modelo de
Programação Linear para o problema de mix de produção da Rentex, e os
primeiros resultados computacionais com análise preliminar.

A formulação foi enriquecida em relação à concepção inicial para capturar o
elemento central do problema: o \textbf{trade-off entre maximizar margem de
contribuição no curto prazo e manter parcerias de longo prazo com clientes
contratuais de menor margem}. Esse elemento distingue o problema da Rentex de
um simples problema de alocação, tornando o uso de Programação Linear
genuinamente necessário para derivar a política ótima.

\newpage

\subsection{Conexão com o Artigo de Referência}

O artigo de \textcite{sobreiro2013}, ``Proposta de uma heurística construtiva
baseada na TOC para definição de mix de produção'', publicado na revista
\textit{Produção} (Qualis A2), é diretamente aplicável ao problema da Rentex.
Esta seção explicita as conexões e o que da metodologia do artigo será utilizado.

\subsubsection{Paralelo entre o Artigo e o Problema Rentex}

O problema tratado por \textcite{sobreiro2013} é, na sua forma geral:

\begin{align}
  &\text{Maximizar} \quad Z = \sum_{j=1}^{n} (P_j - c_j)\, x_j \nonumber \\
  &\text{s.a.} \quad \sum_{j=1}^{n} a_{ij}\, x_j \leq b_i, \quad i = 1, \ldots, m \nonumber \\
  &\phantom{\text{s.a.}} \quad x_j \leq d_j, \quad j = 1, \ldots, n \nonumber
\end{align}

A estrutura é idêntica ao caso Rentex: maximizar a soma das margens de contribuição
sujeita a restrições de capacidade dos recursos e limites de demanda. O problema
da Rentex adiciona restrições de demanda mínima contratual e variáveis de
penalidade por subatendimento, conforme detalhado na seção de formulação.

\subsubsection{O que o Artigo Usa e o que Este Trabalho Adota}

\textcite{sobreiro2013} propõem uma \textbf{heurística construtiva} baseada em
dois elementos: (i) a Teoria das Restrições (TOC), para identificar o recurso
gargalo e ordenar produtos por eficiência no gargalo; e (ii) o problema da
mochila (\textit{knapsack}), para alocar a capacidade disponível de forma gulosa.
O artigo compara sua heurística com a \textbf{solução ótima da Programação Linear
Inteira (PLI)}, que é adotada como referência (\textit{benchmark}).

Este projeto adota uma abordagem diferente e complementar:

\begin{enumerate}[label=\textbf{(\roman*)}]
  \item \textbf{Identificação do gargalo via TOC} (do artigo): a análise da
        eficiência por hora de tear de cada produto -- análoga ao índice
        $R_i = \mathrm{CM}_i / t_{i,\mathrm{BN}_k}$ do artigo -- é usada para
        interpretar a solução ótima e explicar por que certos produtos são
        priorizados.

  \item \textbf{PL como solução exata} (em vez da heurística): como o problema
        tem apenas 4 produtos e 5 recursos, a PL encontra o ótimo global em
        fração de segundo, tornando desnecessária a heurística do artigo. A PL
        é exatamente o \textit{benchmark} que o artigo usa para validar suas
        heurísticas -- estamos operando no patamar de referência máxima.

  \item \textbf{Extensão do modelo com penalidade de relacionamento}: o artigo
        não modela o valor de parcerias de longo prazo. Este projeto incorpora
        variáveis de shortfall e penalidades na função objetivo, o que torna o
        problema não trivial e responde à pergunta gerencial central: \textit{quanto
        vale produzir produtos menos rentáveis para manter parcerias de longo prazo?}

  \item \textbf{Análise de sensibilidade paramétrica}: o artigo sugere como
        trabalho futuro a aplicação das heurísticas em situações práticas e a
        análise de sensibilidade. Este trabalho realiza exatamente isso,
        quantificando o threshold de penalidade a partir do qual o modelo
        recomenda honrar o nível preferido dos clientes contratuais.
\end{enumerate}

\newpage

\subsection{Revisão Bibliográfica}

\subsubsection{O Problema de Mix de Produção com Restrições de Relacionamento}

O problema de definição do mix de produção consiste em determinar as quantidades
a fabricar de cada produto em um horizonte fixo, maximizando a margem de
contribuição total, respeitando as capacidades dos recursos e os limites de
demanda \cite{hillier2006, arenales2007}. Quando existe demanda mínima contratual
e penalidade por subatendimento, o problema passa a capturar o valor das
relações comerciais de longo prazo, tornando-o mais rico e gerencialmente relevante.

\textcite{sobreiro2013} demonstram que, para instâncias de pequeno porte (até
8 produtos e 4 gargalos), a Programação Linear Inteira encontra a solução ótima
em menos de 1 segundo, o que confirma a adequação da PL para o problema da Rentex.

\subsubsection{Programação Linear como Técnica Principal}

A Programação Linear é selecionada por três razões:

\begin{enumerate}[label=\textbf{(\alph*)}]
  \item \textbf{Optimalidade garantida:} o método Simplex encontra o ótimo global
        para qualquer problema LP, ao contrário de heurísticas \cite{hillier2006}.

  \item \textbf{Análise de sensibilidade nativa:} preços-sombra e intervalos de
        ranging emergem diretamente da solução ótima, fornecendo informação
        gerencial sobre o valor marginal de capacidade adicional e sobre a
        robustez do plano a variações nos parâmetros \cite{winston2002}.

  \item \textbf{Escalabilidade para o modelo aprimorado:} a extensão com
        variáveis de shortfall e penalidades mantém a estrutura LP, sem necessitar
        de PLIM, graças à linearidade da penalidade.
\end{enumerate}

\subsubsection{Teoria das Restrições e Gargalo}

A TOC \cite{goldratt2006} identifica que o desempenho do sistema é determinado
por seu gargalo. No contexto da PL, o gargalo é o recurso com maior preço-sombra.
\textcite{sobreiro2013} formalizam a eficiência de cada produto em relação ao
gargalo como $e_j = \mathrm{MC}_j / a_{\text{garg},j}$, que é exatamente o
critério usado neste trabalho para interpretar a priorização dos produtos.

\subsubsection{Modelagem do Valor de Parcerias de Longo Prazo}

A literatura de gestão de operações reconhece que em mercados B2B a relação com
clientes de longo prazo tem valor que vai além da margem de contribuição imediata.
\textcite{arenales2007} discutem como restrições de compromisso mínimo devem
ser incorporadas a modelos de PL para planejamento da produção em empresas com
contratos de fornecimento estável. A abordagem de penalidade por shortfall
adotada neste trabalho transforma as restrições de mínimo contratual de
\textit{hard constraints} em soft constraints parametrizadas, permitindo que o
modelo responda à pergunta: a que custo de margem de contribuição imediata
vale honrar o nível preferido pelo cliente?

\newpage

\subsection{Coleta e Validação dos Dados}

Os dados foram coletados em levantamento conduzido junto ao PCP da Rentex em
abril de 2026, a partir de: roteiros de fabricação do sistema de gestão, relatórios
de capacidade dos últimos doze meses, fichas de custo de produto e histórico de
pedidos dos últimos seis meses.

\subsubsection{Capacidades Semanais Disponíveis}

\begin{table}[H]
\centering
\caption{Capacidades semanais disponíveis validadas}
\label{tab:cap_r2}
\begin{tabular}{clrrc}
\toprule
\textbf{$i$} & \textbf{Recurso} & \textbf{Bruto} & \textbf{Eficiência} & \textbf{$b_i$} \\
\midrule
1 & Teares (h/semana)              & 448 h   & 84,8\% & \textbf{380 h}   \\
2 & Fio de algodão (kg/semana)     & 400 kg  & 100\%  & \textbf{400 kg}  \\
3 & Fio de poliéster (kg/semana)   & 500 kg  & 100\%  & \textbf{500 kg}  \\
4 & Elastano (kg/semana)           &  12 kg  & 100\%  & \textbf{12 kg}   \\
5 & Mão de obra direta (h/semana)  & 1.440 h & 83,3\% & \textbf{1.200 h} \\
\bottomrule
\multicolumn{5}{l}{\footnotesize Bruto dos teares = 5 teares $\times$ 2 turnos $\times$ 8 h $\times$ 5,6 dias úteis médios.}
\end{tabular}
\end{table}

\subsubsection{Produtos, Margens e Níveis de Demanda}

\begin{table}[H]
\centering
\caption{Margens de contribuição e níveis de demanda semanal}
\label{tab:prod_r2}
\begin{tabular}{llrrrrrr}
\toprule
\textbf{$j$} & \textbf{Produto} & \textbf{Preço} & \textbf{CV} & \textbf{MC$_j$} & \textbf{$d^{\min}$} & \textbf{$d^{\text{pref}}$} & \textbf{$d^{\max}$} \\
\midrule
1 & Sarja Barcelos  & 1.800 & 1.100 & \textbf{700}  &   0 & ---  & 200 \\
2 & Sarja Joy       & 2.400 & 1.400 & \textbf{1.000}&   0 & ---  &  80 \\
3 & Forro Padrão    &   600 &   380 & \textbf{220}  &  60 & \textbf{130} & 160 \\
4 & Forro Premium   &   900 &   550 & \textbf{350}  &  20 & \textbf{45}  &  60 \\
\bottomrule
\multicolumn{8}{l}{\footnotesize $d^{\min}$: contrato mínimo inegociável. $d^{\text{pref}}$: nível esperado pelo cliente contratual.}\\
\multicolumn{8}{l}{\footnotesize CV e MC em R\$/lote de 100~m lineares.}
\end{tabular}
\end{table}

\subsubsection{Coeficientes Técnicos}

\begin{table}[H]
\centering
\caption{Coeficientes técnicos $a_{ij}$ (consumo de recurso por lote)}
\label{tab:coef_r2}
\renewcommand{\arraystretch}{1.2}
\begin{tabular}{lcccc}
\toprule
\textbf{Recurso $i$} & \textbf{S. Barcelos} & \textbf{S. Joy} & \textbf{F. Padrão} & \textbf{F. Premium} \\
\midrule
Teares (h/lote) incl. setup & 0,85 & 1,40 & 0,50 & 0,70 \\
Fio algodão (kg/lote)       & 1,35 & 0,95 & 0,00 & 0,00 \\
Fio poliéster (kg/lote)     & 0,00 & 0,35 & 1,60 & 1,55 \\
Elastano (kg/lote)          & 0,00 & 0,07 & 0,00 & 0,08 \\
MOD (h/lote)                & 0,40 & 0,65 & 0,25 & 0,35 \\
\bottomrule
\end{tabular}
\end{table}

O coeficiente de teares incorpora o tempo de setup amortizado pelo tamanho mínimo
de lote contratual (50~m para a Barcelos, 80~m para a Joy), conforme descrito
no Relatório 1.

\newpage

\subsection{Formulação Matemática Final}
\label{sec:form_final}

\subsubsection{Por que o Modelo Básico é Insuficiente}

Um modelo com apenas demanda mínima contratual como \textit{hard constraint}
($x_j \geq d_j^{\min}$) produz um resultado trivialmente previsível: o Simplex
alocaria toda a capacidade de tear aos produtos de maior eficiência (sarjas) e
produziria exatamente o mínimo contratual dos forros. Essa solução dispensa o uso
de um modelo de otimização, pois pode ser obtida por inspeção direta da tabela de
eficiências.

O elemento que torna o problema não trivial é a existência de um \textbf{nível
preferido} pelo cliente -- abaixo do qual existe risco de perda da parceria --
que está acima do mínimo contratual legal. A Rentex precisa decidir: quanto de
margem de contribuição imediata vale sacrificar para honrar o nível preferido e
preservar a relação de longo prazo? Essa é uma decisão de trade-off contínuo
que a Programação Linear é capaz de formalizar e resolver.

\subsubsection{Variáveis de Decisão}

\begin{description}[leftmargin=3cm]
  \item[$x_j \geq 0$] quantidade de lotes do produto $j$ a produzir por semana.
  \item[$s_j^- \geq 0$] \textit{shortfall}: lotes entregues abaixo do nível preferido
        pelo cliente do produto $j$ (definido somente para os produtos com contrato,
        $j \in \{3, 4\}$).
\end{description}

\subsubsection{Parâmetros Adicionais}

\begin{description}[leftmargin=3cm]
  \item[$d_j^{\text{pref}}$] nível de produção preferido pelo cliente contratual
        do produto $j$. Abaixo desse nível, o cliente percebe subatendimento.
  \item[$\rho_j$] penalidade por lote de shortfall (R\$/lote): custo esperado de
        relacionamento por cada lote produzido abaixo de $d_j^{\text{pref}}$.
        Representa o valor semanal amortizado do risco de perda da parceria.
\end{description}

A penalidade $\rho_j$ é calculada como:
\begin{equation}
  \rho_j = \frac{V_j \cdot \pi_j}{52 \cdot \overline{s}_j}
  \label{eq:rho}
\end{equation}
onde $V_j$ é o valor anual do contrato, $\pi_j$ é a probabilidade estimada de
perda do cliente por semana de subatendimento, e $\overline{s}_j$ é o shortfall
esperado por semana de subatendimento. O parâmetro $\rho_j$ é tratado como
variável nos cenários de análise de sensibilidade.

\subsubsection{Modelo de Programação Linear Aprimorado}

\textbf{Função objetivo:}
\begin{equation}
  \max \quad Z = \underbrace{\sum_{j=1}^{4} \mathrm{MC}_j \, x_j}_{\text{MC total bruta}}
                - \underbrace{\sum_{j \in \{3,4\}} \rho_j \, s_j^-}_{\text{custo de relacionamento}}
  \label{eq:fo_r2}
\end{equation}

\textbf{Restrições de capacidade:}
\begin{align}
  0{,}85\,x_1 + 1{,}40\,x_2 + 0{,}50\,x_3 + 0{,}70\,x_4 &\leq 380
    \tag{R1 -- Teares} \label{eq:tear}\\
  1{,}35\,x_1 + 0{,}95\,x_2 &\leq 400
    \tag{R2 -- Fio algodão} \\
  0{,}35\,x_2 + 1{,}60\,x_3 + 1{,}55\,x_4 &\leq 500
    \tag{R3 -- Fio poliéster} \\
  0{,}07\,x_2 + 0{,}08\,x_4 &\leq 12
    \tag{R4 -- Elastano} \\
  0{,}40\,x_1 + 0{,}65\,x_2 + 0{,}25\,x_3 + 0{,}35\,x_4 &\leq 1200
    \tag{R5 -- MOD}
\end{align}

\textbf{Restrições de relacionamento (soft constraints):}
\begin{equation}
  x_j + s_j^- \geq d_j^{\text{pref}}, \quad j \in \{3, 4\}
  \label{eq:soft}
\end{equation}

\textbf{Restrições de demanda:}
\begin{equation}
  x_j \geq d_j^{\min}, \quad x_j \leq d_j^{\max}, \quad \forall j
  \label{eq:dem}
\end{equation}

\textbf{Não negatividade:}
\begin{equation}
  x_j \geq 0, \quad s_j^- \geq 0, \quad \forall j
\end{equation}

O modelo tem \textbf{6 variáveis de decisão} (4 de produção + 2 de shortfall),
\textbf{5 restrições de capacidade}, \textbf{2 soft constraints} de relacionamento
e \textbf{8 restrições de demanda}, totalizando 15 restrições funcionais.

\subsubsection{Verificação de Não Trivialidade}

O plano que atende toda a demanda máxima viola a restrição de teares:
\[
  0{,}85(200) + 1{,}40(80) + 0{,}50(160) + 0{,}70(60)
  = 170 + 112 + 80 + 42 = \mathbf{404 \; \text{h}} > 380 \; \text{h} \quad \times
\]
O modelo não tem solução trivial. Além disso, o nível preferido dos forros
($d_3^{\text{pref}} = 130$, $d_4^{\text{pref}} = 45$) gera um trade-off real:
honrar esses níveis requer realocar horas de tear que de outra forma seriam usadas
para aumentar a produção dos forros ou das sarjas, conforme detalhado na análise.

\newpage

\subsection{Implementação Computacional}
\label{sec:impl}

\begin{lstlisting}[caption={Modelo PL aprimorado da Rentex -- \texttt{rentex\_pl.py}},
                   label={lst:rentex}]
"""
PRO3341 - Projeto Rentex: Mix de Producao com Trade-off de Parcerias
"""
import pulp

produtos   = ['Sarja_Barcelos','Sarja_Joy','Forro_Padrao','Forro_Premium']
recursos   = ['Teares','Fio_Algodao','Fio_Poliester','Elastano','MOD']
contratos  = ['Forro_Padrao','Forro_Premium']

MC     = {'Sarja_Barcelos':700,'Sarja_Joy':1000,'Forro_Padrao':220,'Forro_Premium':350}
b      = {'Teares':380,'Fio_Algodao':400,'Fio_Poliester':500,'Elastano':12,'MOD':1200}
a      = {
    'Teares':       {'Sarja_Barcelos':0.85,'Sarja_Joy':1.40,
                     'Forro_Padrao':0.50,'Forro_Premium':0.70},
    'Fio_Algodao':  {'Sarja_Barcelos':1.35,'Sarja_Joy':0.95,
                     'Forro_Padrao':0.00,'Forro_Premium':0.00},
    'Fio_Poliester':{'Sarja_Barcelos':0.00,'Sarja_Joy':0.35,
                     'Forro_Padrao':1.60,'Forro_Premium':1.55},
    'Elastano':     {'Sarja_Barcelos':0.00,'Sarja_Joy':0.07,
                     'Forro_Padrao':0.00,'Forro_Premium':0.08},
    'MOD':          {'Sarja_Barcelos':0.40,'Sarja_Joy':0.65,
                     'Forro_Padrao':0.25,'Forro_Premium':0.35},
}
d_min  = {'Sarja_Barcelos':0,'Sarja_Joy':0,'Forro_Padrao':60,'Forro_Premium':20}
d_max  = {'Sarja_Barcelos':200,'Sarja_Joy':80,'Forro_Padrao':160,'Forro_Premium':60}
d_pref = {'Forro_Padrao':130,'Forro_Premium':45}

def resolver(rho, b_ov=None):
    _b = {**b, **(b_ov or {})}
    prob = pulp.LpProblem("Rentex", pulp.LpMaximize)
    x = {j: pulp.LpVariable(f"x_{j}", lowBound=0) for j in produtos}
    s = {j: pulp.LpVariable(f"s_{j}", lowBound=0) for j in contratos}

    # Funcao objetivo: MC bruta menos custo de relacionamento
    prob += (pulp.lpSum(MC[j]*x[j] for j in produtos)
             - pulp.lpSum(rho[j]*s[j] for j in contratos))

    for i in recursos:
        prob += pulp.lpSum(a[i][j]*x[j] for j in produtos) <= _b[i], f"Cap_{i}"

    for j in produtos:
        prob += x[j] >= d_min[j], f"Min_{j}"
        prob += x[j] <= d_max[j], f"Max_{j}"

    # Soft constraints: modelo decide se honra o nivel preferido
    for j in contratos:
        prob += x[j] + s[j] >= d_pref[j], f"Pref_{j}"

    prob.solve(pulp.PULP_CBC_CMD(msg=0))
    return prob, x, s

if __name__ == '__main__':
    rho = {'Forro_Padrao': 200, 'Forro_Premium': 300}
    prob, x, s = resolver(rho)
    print(f"Z* = R$ {pulp.value(prob.objective):,.2f}")
    for j in produtos:
        print(f"  {j}: {pulp.value(x[j]):.1f} lotes")
    for j in contratos:
        print(f"  shortfall {j}: {pulp.value(s[j]):.1f} lotes")
\end{lstlisting}

\newpage
